\documentclass[10pt,a4paper]{amsart}
\usepackage[margin=19mm]{geometry}
\usepackage{amsmath,amssymb,mathtools}
\usepackage[scr=rsfs,cal=euler]{mathalfa}
\usepackage{xfrac}
\usepackage[unicode,hidelinks,bookmarks=false]{hyperref}
\newcommand{\ie}{i.e.\ }
\newcommand{\cf}{cf.\ }
\newcommand{\resp}{respectively\ }
\newcommand{\Subsection}{\subsection}
\providecommand{\nobreakdash}{\nobreak\hbox{-}}
\setlength{\emergencystretch}{2em}
\multlinegap0pt
\usepackage{amssymb}
\usepackage{xcolor}
\usepackage{tcolorbox}
\newtheorem{lemma}{Lemma}
\newcommand{\C}{\mathbb{C}}
\renewcommand{\P}{\mathbb{P}}
\newcommand{\Z}{\mathbb{Z}}
\title{The simplicity of the Hodge Bundle}
\author{Anand Patel}
\date{Source: arXiv:2603.19052v3 (CC BY 4.0)}
\begin{document}
\maketitle

\noindent\emph{Source and licence.} The simplicity of the Hodge Bundle. Version arXiv:2603.19052v3 (CC BY 4.0), distributed by arXiv under the Creative Commons Attribution 4.0 International licence, http://creativecommons.org/licenses/by/4.0/, as recorded in the arXiv licence metadata for this exact version and shown on the abstract page https://arxiv.org/abs/2603.19052v3. Full licence notice: \texttt{licenses/arxiv-2603.19052-CC-BY-4.0.txt}.\par\smallskip

\noindent\textit{Editorial scope.} Counted unit: the article's lemma lemma:triples (source0.tex lines 347--350). The article's complete original proof of it is delivered with it (source0.tex lines 352--400, one \texttt{proof} environment of 49 lines). No mathematics is written, repaired or re-derived: the statement and the proof are the article's own lines, in the article's own notation, and no retained line is substituted or re-flowed.  No further source-local context is needed: every cross-reference of every retained line resolves inside the two delivered spans.  Nothing was omitted from any retained span, so the package carries no omission record.  No retained line contains a citation, so the wrapper carries no bibliography and no bibliography entry.  No retained line contains a figure environment, an image-inclusion command or drawing code, so no image is delivered and the package declares no asset dependency.  Editorial formatting only, in addition: an AMS \texttt{amsart} preamble that restates the article's own declarations and macros used by the retained lines, namely the environments \texttt{lemma} on separate counters, together with the standard packages those lines need.  The retained environments print in this document as Lemma 1 in order of appearance, whereas the article prints the same items with the numbering of its own full document.  The complete native source of the article as distributed in the arXiv e-print for this exact version is bundled unchanged as \texttt{sources/196751/source0.tex}.
\begin{lemma}
  \label{lemma:triples} A triple $(m_1, m_2, m_3)$ lies in $T$ if and only if at
  most one of the $m_i$ is $0$ and $m_1 + m_2 + m_3 = g+3$. 
\end{lemma}

\begin{proof}
  Given a triple $(m_1,m_2,m_3) \in T$, we choose generic, mutually coprime,
  square-free polynomials $A,B,C \in \C[x]$ of degrees $m_1, m_2, m_3$,
  respectively. Let $K = \C(x)$.  We define a smooth projective curve $C$ as the
  smooth model of the function field $L = K(y,z)$ where $y^{2} = AC$ and $z^{2}
  = BC$.  Because at most one $m_i = 0$, the polynomials $AC$ and $BC$ cannot
  both be constants; thus the extension $L/K$ forms a fully generated degree $4$
  Galois extension with Galois group $G \cong (\Z/2\Z)^{2}$, guaranteeing $C$ is
  a connected curve.  Since $m_i \equiv g+1 \pmod 2$, the polynomials $AC$ and
  $BC$ have even degrees, implying that $L/K$ is unramified over $x = \infty \in
  \P^{1}$. 

  By construction, the ramification of the cover $C \to \P^{1}$ lies entirely
  over the roots of $A,B$ and $C$.  Over a root $x=p$ of $A$ or $B$, the ramification
  profile is of type $(2,2)$ -- that is, there are two points lying over $p$
  with simple ramification occurring at each point. 

  Over a root $x=c$ of $C(x)$, the quadratic subextension $K(y/z)/K$ evaluates locally to
  $(y/z)^{2} = A(c)/B(c) \neq 0$ and therefore by adjoining $y$ or $z$ we again
  obtain two points of simple ramification over $q$.  

  By the Riemann-Hurwitz formula applied to $L/K$, the genus $g_C$ satisfies: 
  \[2g_{C}-2 = 4(-2) + 2m_1 + 2m_2 + 2 m_3 = -8 + 2(g+3) = 2g-2,\] which
  confirms $g_{C}=g$. The automorphisms $\tau_1, \tau_2$ acting by $(y,z) \mapsto (-y,z)$
 and $(y,z) \mapsto (y,-z)$, respectively, generate the Galois group $G$. Because $L
 = K(z)(y)$ and $y^{2} = AC$, ramification in $C \to C_{2}$ occurs exactly where
 the local valuation of $AC$ in $K(z)$ is odd.  Over a root $x=a$ of $A$, $AC$ has a
 simple zero, and since $C_{2} \to \P^{1}$ is unramified there the valuation
 remains $1$ at the two points lying in $C_{2}$ over $a$.  This gives
 $2m_1$ fixed points so far for $\tau_1$.  

 At the unique point $t \in C_{2}$ lying over a root $x = c$ of $C(x)$, the local function $z$ is a local coordinate. 
 Thus $AC$ vanishes to order $2$ at $t$. Over $t$ the double cover $C \to C_2$ is unramified, and so $\tau_{1}$ swaps the two 
 resulting pre-image points of $t$ in $C$.   Therefore, the fixed point set of $\tau_1$ consists of $2m_1$ points exactly. 

 Similar analyses for $\tau_2$ and $\tau_3 = \tau_1 \tau_{2}$, yield the conclusion that the number of fixed points
 of $\tau_i$ is $2m_i$.


 Conversely, if $C$ is a genus $g$ curve with commuting involutions $\tau_1, \tau_2,
 \tau_3 = \tau_1 \tau_2$ with disjoint fixed point sets of size $2m_{i}$
 respectively, and such that the quotient $C/\langle \tau_1, \tau_2 \rangle
 \simeq \P^{1}$, then the Riemann-Hurwitz analysis performed earlier for the
 cover $C \to \P^{1}$ shows that
 $m_1 + m_2 + m_3 = g+3$ holds. Furthermore, each $m_{i} \equiv g+1 \pmod 2$ by
 applying Riemann-Hurwitz to the quotients $C \to C /\tau_i$. Two of the $m_i$, say $m_1, m_2$ cannot simultaneously be $0$ since then $C \to
 \P^{1}$ would be a degree $4$ connected unramified
 cover of $\P^{1}$, which is impossible. The lemma is shown. 
\end{proof}

\end{document}
